% TOP_PROBLEM: {"id": 9400090, "title": "Szpiro's conjecture", "category_id": 1, "category": {"id": 1, "name": "number_theory", "display_name": "Number Theory", "description": "Properties of integers, prime numbers, Diophantine equations.", "slug": "number-theory", "order_index": 1, "created_at": "2026-07-31T15:26:25.670Z"}, "status": "partially_solved"}
% FIRST_POSTED: 2026-09-25
% ENTRY_KIND: statement_only
% !TeX program = lualatex
% Definition and sources only; this file is not an LLM solution attempt.
\documentclass[11pt]{article}
\usepackage[margin=25mm]{geometry}
\usepackage{fontspec}
\setmainfont{Latin Modern Roman}
\usepackage{amsmath,amssymb,mathtools}
\usepackage{unicode-math}
\setmathfont{Latin Modern Math}
\usepackage{xurl,hyperref}
\hypersetup{colorlinks=true,urlcolor=blue}
\setcounter{secnumdepth}{0}
\setlength{\parindent}{0pt}
\setlength{\parskip}{5pt}
\setlength{\emergencystretch}{3em}
\begin{document}
\section*{\texorpdfstring{Szpiro's conjecture}{Szpiro's conjecture}}\label{9400090}
\subsection{Definitions and mathematical statement}
For an elliptic curve $E/{\mathbb{Q}}$, $\Delta_E\in{\mathbb{Z}}\setminus\{0\}$ is the discriminant of a global minimal Weierstrass equation. Its conductor $N_E=\prod_p p^{f_p}$ records the local conductor exponents of its Tate-module representation. The conjecture is
\[
 \forall\varepsilon>0\ \exists C_\varepsilon>0\ \forall E/{\mathbb{Q}}:\qquad
 |\Delta_E|\le C_\varepsilon N_E^{6+\varepsilon}.
\]
The constant is independent of the curve. The selected exponent is the classical $6+\varepsilon$ conductor--discriminant bound, not a differently normalized modified Szpiro assertion.
\par\smallskip\textit{Formulation and concept references:} \href{https://arxiv.org/html/2308.07114}{[S1]}.

\subsection{Short English statement}
Can every elliptic curve's minimal discriminant be bounded uniformly by essentially the sixth power of its conductor?
\subsection{Sources}
\begin{itemize}
\item[S1]H. Pasten, Szpiro's conjecture when the denominator of the j-invariant is small, arXiv v2 (15 August 2023).
\url{https://arxiv.org/html/2308.07114}.
\textit{Formulation source.}
\end{itemize}
\end{document}
